\chapter{Intersection Multiplicities and Homological Bounds}
\label{ch:intersection-multiplicities}

At an isolated intersection the ordinary tensor product records the
intersection scheme. Higher Tor groups measure the failure of tensor
product to preserve a resolution. Taking their alternating length makes
the resulting number additive. In a power series ring, the diagonal
presents this alternating length as a Koszul Euler characteristic.
Hilbert--Samuel theory then supplies its sign.
\cite[Chapter V]{serre-local}

\section{The Tor Euler characteristic}

Write \(\ell_R\) for module length. All complexes in this chapter, except
for a Cech complex written cohomologically, have homological grading.
The total differential on \(P_i\otimes Q_j\) is
\(d_P\otimes 1+(-1)^i1\otimes d_Q\).

\begin{definition}[The isolated-intersection Euler characteristic]\label{res:I01}
Let \((R,\mathfrak m)\) be a regular local ring of dimension \(d\).
For finite \(R\)-modules \(M,N\) with \(M\otimes_R N\) of finite length, set
\[
 \chi_R(M,N)=\sum_{i=0}^d(-1)^i
                \ell_R\bigl(\operatorname{Tor}_i^R(M,N)\bigr).
\]
The terms have finite length, all higher terms vanish, and the integer is
symmetric in \(M,N\).
\end{definition}
\begin{proof}
A finite module has a resolution by finite free modules of length at
most \(d\) by \cite[CM-12 or CM-19]{core}. The Tor modules are finite
and the resolution localizes. A localized Tor module vanishes when
one localized argument vanishes. Conversely, two nonzero finite
modules over \(R_{\mathfrak p}\) have nonzero tensor product because reduction
to its residue field gives the tensor product of two nonzero vector
spaces by Nakayama's lemma. Thus
\[
 \operatorname{Supp}(M\otimes_R N)
   =\operatorname{Supp}(M)\cap\operatorname{Supp}(N).
\]
This support is contained in \(\{\mathfrak m\}\) under the hypothesis.
A finite module with this support is killed by some power of
\(\mathfrak m\), and its successive \(\mathfrak m\)-adic quotients are
finite-dimensional over the residue field. It has finite length.
For symmetry, use the total tensor product of finite free resolutions
of \(M,N\). Taking homology in the two possible directions gives the
two computations of Tor, and interchange of the factors sends
\(p\otimes q\) to \((-1)^{ij}q\otimes p\) in bidegree \((i,j)\).
\cite[\S2]{dugger}
\end{proof}

\begin{proposition}[Additivity and reduction to prime quotients]\label{res:I02}
For a fixed finite \(N\), the Euler characteristic is additive on an
exact sequence of finite modules whose supports meet the support of
\(N\) only at \(\mathfrak m\). The same holds in the other argument.
For prime filtrations of \(M\) and \(N\) with factors
\(R/\mathfrak p_a\) and \(R/\mathfrak q_b\), respectively, one has
\[
 \chi_R(M,N)=
 \sum_{a,b}\chi_R(R/\mathfrak p_a,R/\mathfrak q_b).
\]
Each pair on the right still satisfies the isolated-support condition.
The maxima of the dimensions of the respective factors are
\(\dim M\) and \(\dim N\) for nonzero modules.
\end{proposition}
\begin{proof}
The long exact Tor sequence is finite and consists of finite-length
modules by \ref{res:I01}. Its alternating length is zero, since the
length of a term is the sum of the lengths of the incoming and outgoing
images. Regrouping yields additivity.

In a nonzero finite module choose a nonzero element with maximal
annihilator among the annihilators of nonzero elements. Such an
annihilator exists by the ascending chain condition on ideals. If
\(ab\) kills this element but \(b\) does not, multiplication by \(b\)
produces a nonzero element whose annihilator contains the original
annihilator and \(a\). Maximality puts \(a\) in the original
annihilator, so it is prime. Remove the cyclic submodule and repeat
in the quotient. The preimages form an ascending chain of submodules,
which terminates by Noetherianness. Localization of each short exact
sequence shows that the support is the union of the supports of the
factors. This proves the support and dimension claims, and iterated
additivity gives the formula. Equivalently the Euler pairing descends
to the Grothendieck groups of finite modules on the two fixed closed
supports with isolated intersection, where the relations impose
additivity on short exact sequences. \cite[\S2, Lemma 2.2]{dugger}
\end{proof}

\begin{proposition}[Passage to the completion]\label{res:I03}
For an isolated intersection in a regular local ring,
\[
 \chi_{\widehat R}(\widehat M,\widehat N)=\chi_R(M,N).
\]
The dimensions of the nonzero modules and the regular local dimension
of \(R\) are preserved by maximal-ideal completion.
\end{proposition}
\begin{proof}
The ring map to the completion is flat and induces an isomorphism on
residue fields by \cite[ALG-05-C04-S01]{p103}. Tensoring a finite free
resolution and then commuting flat base change with homology identifies
each completed Tor module with the corresponding Tor over
\(\widehat R\). A composition series of a finite-length module stays a
composition series after base change, since each simple factor becomes
the same residue field. In particular the tensor product still has
finite length. The lengths of all maximal-ideal power quotients are
unchanged, so the Hilbert--Samuel dimension formula
\cite[ALG-05-C03-S02]{p103} preserves dimension. For the ring, the
dimension of the maximal ideal modulo its square also stays the same.
The definition of regular local ring makes \(\widehat R\) regular.
\end{proof}

\section{Koszul Euler characteristics}

\begin{definition}[Multiplicity in a specified degree]\label{res:I04}
Let \((A,\mathfrak n)\) be Noetherian local and let \(E\ne 0\) be a
finite \(A\)-module. Suppose \(J\subseteq\mathfrak n\) and \(E/JE\)
has finite length. Write \(P_{J,E}(t)\) for the eventual polynomial
\(\ell_A(E/J^{t+1}E)\) for integer \(t\ge 0\). For an integer
\(r\ge\dim_A E\) let
\[
 e^{(r)}_J(E)=r!\,[t^r]P_{J,E}(t).
\]
Its value at \(r=\dim_A E\) is the Hilbert--Samuel multiplicity
\(e_J(E)\). For the zero module put \(e^{(r)}_J(0)=0\).
\end{definition}

Existence, degree \(\dim_A E\), and positivity of the leading
coefficient for a nonzero module follow from
\cite[ALG-05-C03-S02]{p103}. The notation \(e^{(r)}\) retains the
degree in which the coefficient is taken. A module of smaller
dimension contributes zero in that degree.

\begin{theorem}[Serre's Koszul Euler formula]\label{res:I05}
Let \(z=(z_1,\ldots,z_r)\) generate \(J\subseteq\mathfrak n\), where
\(E/JE\) has finite length. Write \(K_\bullet(z;E)\) for the Koszul
complex. For nonzero \(E\) one has \(\dim_A E\le r\), and
\[
 \sum_{i=0}^r(-1)^i\ell_A H_i(K_\bullet(z;E))=e^{(r)}_J(E).
\]
In particular, when \(\dim_A E=r\) this multiplicity is a positive
integer.
\end{theorem}
\begin{proof}
The dimension bound follows by applying the height theorem to the
image of \(J\) in \(A/\operatorname{Ann} E\), where its radical is
the maximal ideal.
\cite[ALG-05-C03-R01 and ALG-05-C03-S01]{p103}
For \(r=0\) the assertion reads \(\ell_A E=\ell_A E\). For the zero
module all lengths are zero, so suppose \(r>0\) and \(E\ne 0\).

Multiplication by each \(z_j\) on the Koszul complex is homotopic to
zero via exterior multiplication by its basis vector. Thus every
homology group vanishes outside \(V(J)\) and is supported in
\(\operatorname{Supp}(E)\). It has finite length, since
\(V(J)\cap\operatorname{Supp}(E)\subseteq\{\mathfrak n\}\).
Write \(K_i=K_i(z;E)\) and \(Z_i=\ker(d_i)\). For an integer \(n>r\)
the submodules \(F^nK_i=J^{n-i}K_i\) form a subcomplex. For each \(i\),
Artin--Rees applied to \(Z_i\subseteq K_i\) gives, for all sufficiently
large \(n\),
\[
 Z_i\cap J^{n-i}K_i
     =J\bigl(Z_i\cap J^{n-i-1}K_i\bigr).
\]
A cycle on the left is then \(\sum_j z_j w_j\) with each \(w_j\)
a cycle in \(J^{n-i-1}K_i\). The Koszul homotopy identifies \(z_jw_j\)
with \(d(e_j\wedge w_j)\), and its preimage belongs to
\(J^{n-i-1}K_{i+1}=F^nK_{i+1}\). In last degree the wedge is zero
and the same identity states that the term is zero. This proves
acyclicity of \(F^nK_\bullet\) simultaneously in all degrees for
large \(n\). The quotient complex therefore has the original
homology.

Each quotient term has finite length. For \(a\ge 1\),
\(E/J^aE\) has support at most the closed point. For a bounded complex
of finite length modules, the alternating sum of the lengths of
the terms equals that of the homology. This follows by applying
length additivity to the image and kernel sequences for every
differential. Using \(K_i=E^{\binom ri}\), the alternating length
of the quotient complex is
\[
 \sum_{i=0}^r(-1)^i\binom ri P_{J,E}(n-i-1).
\]
This is the \(r\)-fold backward difference of a polynomial of degree
at most \(r\) evaluated at \(n-1\). Each difference reduces the
positive degree by one and multiplies the leading coefficient by
the old degree. The expression equals \(r!\,[t^r]P_{J,E}(t)\).
Positivity in degree \(r\) follows from the Hilbert--Samuel dimension
formula and integrality from the length sum. The proof also gives
additivity of the degree-\(r\) formula on finite modules supported
at the closed point after quotienting by \(J\). Apply the long exact
Koszul homology sequence to a short exact sequence of coefficients.
\cite[Chapter IV, \S A.3]{serre-local}
\end{proof}

\begin{corollary}[One complete-intersection factor]\label{res:I06}
If \(z\) is an \(R\)-regular sequence of length \(r\) in the maximal
ideal of a Noetherian local ring \(R\), and \(N\ne 0\) is finite
with \(N/(z)N\) of finite length, the Tor Euler characteristic
of \(R/(z),N\), defined by the sum through degree \(r\), equals
\(e^{(r)}_{(z)}(N)\).
It is zero for \(\dim_R N<r\) and positive for \(\dim_R N=r\).
\end{corollary}
\begin{proof}
The Koszul complex of the regular sequence is a finite free
resolution of \(R/(z)\). \cite[Chapter 5]{core}
After tensoring with \(N\) it is precisely the complex of
\ref{res:I05}. Thus all assertions hold without a characteristic
restriction or regularity assumption on \(R\) itself. When \(R\)
is regular of dimension \(d\), \(R/(z)\) has dimension \(d-r\) by
successive dimension drops, so this is the isolated-intersection
vanishing and proper-intersection positivity criterion for this
particular first argument.
\end{proof}

\section{The completed exterior tensor product over a field}

\begin{proposition}[Construction and separate exactness]\label{res:I07}
Let \(k\) be a field and put
\[
 A=k[[x_1,\ldots,x_a]],\qquad B=k[[y_1,\ldots,y_b]],
 \qquad S=k[[x_1,\ldots,x_a,y_1,\ldots,y_b]] .
\]
For finite \(A\)- and \(B\)-modules \(M,N\) there is a natural finite
\(S\)-module
\[
 M\boxtimes_k N=(S\otimes_A M)\otimes_B N.
\]
It is exact in either variable on finite modules. It has the
description
\[
 M\boxtimes_k N
   \ \cong\
 \varprojlim_{u,v}
   (M/(x)^uM)\otimes_k(N/(y)^vN).
\]
\end{proposition}
\begin{proof}
The ordinary tensor expression is finite over \(S\), for generators
of each factor give generators indexed by their pairs. Its quotient by
\((x)^uS+(y)^vS\) is the displayed finite-level tensor product over \(k\).
Indeed the quotient of \(S\) is
\(A/(x)^u\otimes_k B/(y)^v\) by the monomial bases, and the claimed
quotient follows by the universal balanced tensor maps.
The ideals with \(u=v=n\) are cofinal with powers of \((x,y)S\).
\[
 (x,y)^{2n-1}S\ \subseteq\ (x)^nS+(y)^nS\
             \subseteq\ (x,y)^nS .
\]
Finite modules over the complete Noetherian ring \(S\) are complete
by the completion tensor comparison of
\cite[ALG-05-C04-R03]{p103}. This proves the inverse-limit formula.

For an exact sequence \(0\to M'\to M\to M''\to0\), use the exact
finite-level sequence
\[
 0\to M'/(M'\cap(x)^nM)\to M/(x)^nM\to M''/(x)^n M''\to0 .
\]
Tensor with \(N/(y)^nN\) over the field. The left terms now have
surjective transition maps. A compatible class in the right terms
lifts to a compatible class in the middle. After choosing the
next lift, its discrepancy with the preceding lift comes from a
left term, and that discrepancy lifts to the next left term.
Subtracting it supplies the required correction. Kernel exactness
and injectivity are checked termwise. Artin--Rees makes
\(M'\cap(x)^nM\) cofinal with \((x)^nM'\), and using the respective
cofinal indices gives mutually inverse maps on limits after tensoring
with the truncations of \(N\). The limit sequence is the exterior
tensor sequence, and the same argument interchanges the roles.
\cite[Chapter V, \S B.2]{serre-local}
\end{proof}

\begin{lemma}[Dimensions add over the field]\label{res:I08}
For nonzero finite \(M,N\) in \ref{res:I07}, with
\(\mathfrak s=(x,y)S\), one has
\[
 \dim_S(M\boxtimes_k N)=\dim_A M+\dim_B N,\qquad
 e_{\mathfrak s}(M\boxtimes_k N)=e_{(x)}(M)e_{(y)}(N).
\]
\end{lemma}
\begin{proof}
Let \(T=M\boxtimes_k N\). For any \(j\ge 0\) the natural product map
from the vector spaces of the two associated graded modules identifies
\[
 \mathfrak s^jT/\mathfrak s^{j+1}T
   \ \cong\
 \bigoplus_{u+v=j}
       \frac{(x)^uM}{(x)^{u+1}M}
              \otimes_k\frac{(y)^vN}{(y)^{v+1}N}.
\]
To verify both surjectivity and injectivity, work modulo
\((x)^{j+1}S+(y)^{j+1}S\), which is contained in
\(\mathfrak s^{j+1}\). Split the flags in the two finite-dimensional
truncated vector spaces over \(k\). In their tensor product,
\(\mathfrak s^j\) acts through the sum of the images of
\((x)^uM\otimes (y)^vN\) with \(u+v=j\). These images are the sum of
split tensors with total filtration index at least \(j\).
The quotient removes precisely those with index at least \(j+1\).
The natural map therefore has the displayed description, independent
of the choices of splittings.

For a nonzero finite module \(U\) of dimension \(s\) over a local
ring with the indicated maximal-ideal filtration, let \(h_U(j)\)
be its \(j\)-th layer length. The Hilbert--Samuel dimension formula
gives
\[
 \sum_{j\ge 0}h_U(j)t^j=\frac{Q_U(t)}{(1-t)^s},
 \qquad Q_U(1)=e(U)>0 ,
\]
where \(Q_U(t)\) is a polynomial and \(e(U)\) is its maximal-ideal
multiplicity. To see the asserted value, first sum the cumulative
lengths. An eventual polynomial of degree \(s\) expands in
\(\binom{j}{i}\) with highest coefficient \(e(U)\).
The corresponding series for those polynomials are
\(t^i/(1-t)^{i+1}\), and a finite initial correction is a polynomial.
Multiplication by \(1-t\) gives the series of layer lengths and
preserves the stated numerator at \(1\). This argument includes
\(s=0\), where the layer series is already a polynomial.
The direct-sum calculation multiplies the two layer series. The
numerator at \(1\) is the product of two positive numbers, so the
pole order and numerator give exactly the asserted dimension and
multiplicity. \cite[Chapter V, \S B.3]{serre-local}
\end{proof}

\begin{theorem}[Diagonal reduction in equal characteristic]\label{res:I09}
Let \(R=k[[x_1,\ldots,x_d]]\) and let \(M,N\ne0\) be finite
\(R\)-modules such that \(M\otimes_R N\) has finite length. Regard
\(N\) as a module over the copy with variables \(y\) and set
\(T=M\boxtimes_k N\) and \(J=(x_1-y_1,\ldots,x_d-y_d)\subset S\).
Then
\[
 \operatorname{Tor}_i^R(M,N)\cong H_i(K_\bullet(x-y;T)),\qquad
 \chi_R(M,N)=e^{(d)}_J(T).
\]
In particular \(\dim M+\dim N\le d\), with zero Euler
characteristic for strict inequality and positive Euler
characteristic for equality. The conclusion holds for any
equal-characteristic regular local ring.
\end{theorem}
\begin{proof}
The variable substitution \(z_i=x_i-y_i\) gives the isomorphism
\(S=k[[z_1,\ldots,z_d,y_1,\ldots,y_d]]\), and \(S/J\) identifies
with \(R\). The \(z_i\) form a regular sequence with Koszul
resolution of \(S/J\).

Take bounded finite free resolutions \(P,Q\) over the \(x\) and
\(y\) copies, respectively. Their exterior tensor double complex
has finite free \(S\)-terms. Separate exactness in
\ref{res:I07}, applied successively to the two resolutions, gives
a finite free resolution of \(T\) on taking the total complex.
To justify this computation, filter the total complex in either
direction. Homology in the first direction leaves only the
augmented column or row, and homology in the second leaves only
\(T\) in degree zero. After tensoring this total resolution with
\(S/J\), the terms and differentials identify with the total
complex \(P\otimes_R Q\). One can again take homology first on the
free resolution \(P\) to obtain \(\operatorname{Tor}^R(M,N)\).
Computing \(\operatorname{Tor}^S(S/J,T)\) instead with the
diagonal Koszul resolution gives the homology in the assertion.
The two computations agree, by the same double complex argument.
The \(S\)-action factors through \(R\) and the finite lengths
agree for these two coefficient rings.

In degree zero,
\(T/JT=M\otimes_R N\) has finite length. Thus \ref{res:I05}
applies to \(J\) with \(d\) generators. By \ref{res:I08},
\(\dim T=\dim M+\dim N\), giving the strict and equality cases.
For a general equal-characteristic regular local ring apply
\ref{res:I03} and use the coefficient field and power series
description of its completion from \ref{res:C04}.
\cite[Chapter V, \S B.3, Theorem 1]{serre-local}
\end{proof}

\section{Mixed characteristic and the dimension bound}

\begin{lemma}[Euler cancellation across a regular element]\label{res:I10}
Let \(S\) be Noetherian local, let \(h\) be a nonzero divisor in
its maximal ideal, and set \(A=S/(h)\). Let \(M,N\) be finite
\(A\)-modules such that \(M\otimes_A N\) has finite length.
Suppose \(M\) has a finite free \(A\)-resolution and \(N\) has a
finite free \(S\)-resolution. Then the alternating length of
\(\operatorname{Tor}_i^S(M,N)\) is zero.
\end{lemma}
\begin{proof}
Let \(P\) be a bounded finite free \(A\)-resolution of \(M\) and
\(Q\) a bounded finite free \(S\)-resolution of \(N\). Consider
the bounded double complex \(P\otimes_A(A\otimes_S Q)\).
Its total complex identifies with the total complex of
\(P\otimes_S Q\). In the direction of \(P\), exactness against
each free \(S\)-module \(Q_j\) leaves \(M\otimes_S Q\), so the
abutment is \(\operatorname{Tor}^S(M,N)\).
In the direction of \(Q\), one obtains the groups
\(P_i\otimes_A \operatorname{Tor}_j^S(A,N)\), since each \(P_i\)
is \(A\)-free. A two-term free \(S\)-resolution of \(A\) is
\([S\xrightarrow{h}S]\). Multiplication by \(h\) on \(N\) is
zero. Thus \(\operatorname{Tor}_j^S(A,N)\cong N\) for \(j=0,1\)
as an \(A\)-module, with all other groups zero. For this
identification, multiplication by a scalar in \(S\) on either
argument is represented by multiplication by that scalar on
the tensor complex. Hence the action obtained from \(A\) on
the first argument agrees with that of \(S\) on \(N\) and
factors through \(A\). Taking the next
homology gives \(\operatorname{Tor}_i^A(M,N)\) in each of the
two rows \(j=0,1\). These entries have finite length by the
same support argument as in \ref{res:I01}. Length over \(A\)
equals length over \(S\).

For each succeeding page the total alternating length is
preserved since pairs of boundary terms contribute with opposite
signs, as each differential changes total degree by one.
The double complex is bounded, so the last page gives a finite
filtration on each abutment group. The two identical rows
contribute with opposite signs in total degree, which proves
the cancellation. \cite[Chapter V, \S B.5]{serre-local}
\end{proof}

\begin{lemma}[Diagonal reduction over a discrete valuation ring]\label{res:I11}
Let \(V\) be a complete discrete valuation ring with uniformizer
\(\pi\) and residue field \(k\). Put
\[
 R=V[[x_1,\ldots,x_a]],\qquad
 S=V[[x_1,\ldots,x_a,y_1,\ldots,y_a]] .
\]
Let \(M,N\) be nonzero finite \(R\)-modules on which multiplication
by \(\pi\) is injective. Define \(T=(S\otimes_{V[[x]]}M)
\otimes_{V[[y]]}N\), with \(N\) written in the \(y\)-variables.
Then multiplication by \(\pi\) on \(T\) is injective and
\[
 T/\pi T\cong (M/\pi M)\boxtimes_k(N/\pi N),\qquad
 \dim_S T=\dim_R M+\dim_R N-1 .
\]
For \(J=(x_1-y_1,\ldots,x_a-y_a)\), the diagonal Koszul complex
\(K_\bullet(x-y;T)\) computes \(\operatorname{Tor}^R(M,N)\).
\end{lemma}
\begin{proof}
The rings in the statement are regular local, with respective
systems of parameters \(\pi,x_1,\ldots,x_a\) and
\(\pi,x_1,\ldots,x_a,y_1,\ldots,y_a\). This is the power series
construction of \ref{res:C04}. Take bounded finite free resolutions
\(P,Q\) over the two copies of \(R\). Form the total complex \(D\) of
the finite free double complex
\(S\otimes_{V[[x]]}P\otimes_{V[[y]]}Q\). Its zeroth homology is
\(T\) by right exactness of tensor product in the two arguments.
Since \(\pi\) is injective on \(M,N\), the exact sequences of
complexes for multiplication by \(\pi\) show that \(P/\pi P\)
and \(Q/\pi Q\) are finite free resolutions of \(M/\pi M\)
and \(N/\pi N\). Thus \(D/\pi D\) is exactly their exterior tensor
double complex from the proof of \ref{res:I09}.
It has no positive homology by \ref{res:I07}, and its zeroth
homology is the exterior tensor product in the statement.
The long exact sequence for
\(0\to D\xrightarrow{\pi}D\to D/\pi D\to0\) implies
that multiplication by \(\pi\) acts surjectively on every
\(H_i(D)\) with \(i>0\). Each is finite, hence zero by Nakayama.
The segment at degree zero gives the injectivity on \(T\)
and the claimed quotient identification. In particular \(D\)
resolves \(T\).

For a finite nonzero module \(U\) over a Noetherian local ring
and a nonzero divisor \(u\) on \(U\) in the maximal ideal,
\(\dim(U/uU)=\dim U-1\). Here is the dimension argument.
Each minimal prime of the support of \(U\) avoids \(u\),
since such a prime is an associated prime. Every prime
in \(\operatorname{Supp}(U/uU)\) contains \(u\) and lies
strictly over a minimal prime of \(\operatorname{Supp}(U)\).
Prepending that prime to a longest chain above the given prime
gives the upper bound. Conversely, prepend \(u\) to a system of
parameters on \(U/uU\). The resulting sequence has finite-colength
quotient on \(U\), so the height bound gives
\(\dim U\le 1+\dim(U/uU)\). For the associated prime and parameter
facts used here see \cite[Chapter 6]{core} and
\cite[ALG-05-C03-S01]{p103}.
Apply this drop to \(M,N,T\), with \(u=\pi\). Their special-fiber
dimensions are related by \ref{res:I08}, proving the formula.
To compute Tor, base change \(D\) to \(S/J\cong R\). The total
complex becomes \(P\otimes_R Q\), as in \ref{res:I09}. Compute
the same Tor groups with the Koszul resolution of the diagonal
over \(S\). The substitution \(z_i=x_i-y_i\) is valid over \(V\),
so this is a regular sequence. \cite[Chapter V, \S B.5, Lemma 1]{serre-local}
\end{proof}

\begin{theorem}[Unramified mixed characteristic]\label{res:I12}
Let \(R\) be a regular local ring of mixed characteristic
\((0,p)\) such that \(p\notin\mathfrak m^2\), of dimension
\(d\). For nonzero finite modules meeting only at the closed
point, \(\dim M+\dim N\le d\). The number \(\chi_R(M,N)\)
is zero if this inequality is strict and positive if it is
equality.
\end{theorem}
\begin{proof}
By \ref{res:I03} pass to the completion, where \ref{res:C04}
identifies \(R\) with \(V[[x_1,\ldots,x_{d-1}]]\) for
a complete discrete valuation ring \(V\) of uniformizer \(p\).
By \ref{res:I02} it suffices to prove the claims for prime
quotients. Indeed all pairs of prime-filtration factors
have dimensions bounded by those of the original modules.
If the original dimension sum is \(d\), the factors attaining
the two maxima give at least one strictly positive term, and
all other terms are nonnegative.

A prime quotient is either killed by \(p\) or \(p\) acts
injectively on it. If both arguments have injective \(p\)-action,
apply \ref{res:I11}. The diagonal has \(d-1\) generators and
\(T/JT\cong M\otimes_R N\) has finite length. By \ref{res:I05},
\(\dim T\le d-1\) and the sign follows from the comparison
\(\dim T=\dim M+\dim N-1\).

Suppose exactly one factor, say \(M\), is killed by \(p\).
Write \(\bar R=R/pR\). For a finite free \(R\)-resolution
\(Q\) of \(N\), the homology sequence of multiplication by \(p\)
shows \(Q/pQ\) resolves \(N/pN\) over \(\bar R\). Termwise,
\(M\otimes_R Q=M\otimes_{\bar R}(Q/pQ)\). Thus both the
Tor groups and lengths agree with those for \(M,N/pN\) over
\(\bar R\). The dimensions of \(M\) over \(R\) and \(\bar R\)
agree, and \(\dim(N/pN)=\dim N-1\) by the drop used in
\ref{res:I11}. The ring \(\bar R\) is equal-characteristic
of dimension \(d-1\), so \ref{res:I09} gives the result.

If both arguments are killed by \(p\), apply \ref{res:I09}
over \(\bar R\) to see that their dimension sum is at most
\(d-1\). Their alternating length over \(R\) vanishes by
\ref{res:I10}, with \(S=R\), \(A=\bar R\), \(h=p\).
All modules admit the required bounded resolutions since the
rings are regular. This completes the three cases.
\cite[Chapter V, \S B.5, Theorem 2]{serre-local}
\end{proof}

\begin{theorem}[Serre's dimension bound]\label{res:I13}
Let \(R\) be a regular local ring of dimension \(d\) in any
characteristic. For nonzero finite modules with finite-length
tensor product,
\[
 \dim_R M+\dim_R N\le d.
\]
In particular, if \(\mathfrak r\) is minimal over
\(\mathfrak p+\mathfrak q\) for two prime ideals in a regular
local ring, then
\(\operatorname{ht}\mathfrak r\le
  \operatorname{ht}\mathfrak p+\operatorname{ht}\mathfrak q\).
\end{theorem}
\begin{proof}
The equal-characteristic and unramified cases are
\ref{res:I09} and \ref{res:I12}. By \ref{res:I03} it remains
to consider a complete mixed-characteristic regular local \(R\)
with \(p\in\mathfrak m^2\). The coefficient-ring construction
of \ref{res:C03} supplies a complete discrete valuation ring
\(V\to R\) with uniformizer \(p\) and residue-field isomorphism.
For a minimal system of parameters \(x_1,\ldots,x_d\) of \(R\),
the natural map
\[
 S=V[[X_1,\ldots,X_d]]\longrightarrow R,\qquad X_i\longmapsto x_i
\]
is surjective. At each order the parameters generate the
associated graded ring over the residue field, so one lifts
successive homogeneous pieces using coefficients in \(V\).
The resulting series converges to the prescribed element.
Since \(p\in\mathfrak m^2\), choose \(f\in (X_1,\ldots,X_d)^2\)
mapping to \(p\) by lifting its expression in products \(x_ix_j\).
The kernel contains \(h=p-f\), which is part of a minimal
system of generators of the maximal ideal of \(S\).
The quotient \(S/(h)\) is regular local and a domain of
dimension \(d\) by \cite[ALG-05-C03-S03]{p103}. A nonzero
ideal in this domain has quotient of dimension at most
\(d-1\), since any prime chain in its vanishing locus may be
lengthened below by the zero prime. Therefore the induced
surjection \(S/(h)\to R\) is an isomorphism.

View \(M,N\) as \(S\)-modules. Their dimensions and finite
tensor length are unchanged. By \ref{res:I10},
\(\chi_S(M,N)=0\). The regular local ring \(S\) is unramified
of dimension \(d+1\). By \ref{res:I12} the dimension sum is
at most \(d+1\), and equality would make \(\chi_S(M,N)>0\).
The bound follows.

For the height statement localize at \(\mathfrak r\), which
preserves regularity by \cite[CM-19]{core}. The localized
prime quotients meet only at the closed point since
\(\mathfrak r\) is minimal over their sum. Their dimensions are
\(\operatorname{ht}\mathfrak r-\operatorname{ht}\mathfrak p\)
and \(\operatorname{ht}\mathfrak r-\operatorname{ht}\mathfrak q\),
by the height and dimension formula in a Cohen--Macaulay local
ring \cite[CM-07]{core}. Heights of these smaller primes do
not change upon localization, since all primes under them
are contained in \(\mathfrak r\). Substitution into the
dimension bound gives the stated inequality.
\cite[Chapter V, \S B.6]{serre-local}
\end{proof}

\section{Finite free complexes and examples}

\begin{proposition}[A depth bound for short complexes]\label{res:I14}
Let \((A,\mathfrak n)\) be Noetherian local. Let \(F_\bullet\)
be a finite free complex in degrees \(0,\ldots,s\), and let
\(U\ne0\) be finite. Suppose all \(H_i(F_\bullet\otimes_A U)\)
have finite length and some is nonzero. Then
\(s\ge\operatorname{depth}_A U\). In particular, a nonexact
finite free complex over a Cohen--Macaulay local ring with
finite-length homology has length at least the dimension
of the ring.
\end{proposition}
\begin{proof}
Choose generators of \(\mathfrak n\). Let \(C^\bullet\) be
the augmented cohomological Cech complex from \ref{res:L05}.
Its terms are flat. Give \(C^q\otimes F_i\otimes U\) total
cohomological degree \(q-i\). In the filtration taking
\(F_\bullet\otimes U\) homology first, flatness gives
\(C^\bullet\otimes H_i(F_\bullet\otimes U)\) in the \(i\)-th
row. Each finite-length homology group is killed by a power
of \(\mathfrak n\). Localizing such a group at any generator of
\(\mathfrak n\) makes it zero. Its Cech complex is therefore
the group itself in degree zero, and all surviving terms
have total degree \(-i\le 0\). There is a nonzero term, and
no later differential can leave or enter the only remaining
Cech degree. The abutment has a nonzero group in a
nonpositive degree.

In the other filtration the initial homology groups are
\(H^q_{\mathfrak n}(U)^{\operatorname{rank}F_i}\). They are
zero for \(q<\operatorname{depth}_A U\) by \ref{res:V01}.
If \(\operatorname{depth}_A U>s\), every nonzero term then
has \(q-i>0\). This contradicts the nonpositive-degree
abutment from the first filtration. Both filtrations are
bounded and converge to the homology of the same finite
total complex. Apply the assertion with \(U=A\) for the
last claim. \cite[\S 9.4, Corollary 9.4.3]{bruns}
\end{proof}

\begin{corollary}[Projective dimension and isolated support]\label{res:I15}
Let \(M,N\ne 0\) be finite modules over a Noetherian local
ring \(R\) with \(M\otimes_R N\) of finite length and
\(\operatorname{pd}_R M\) finite. Then
\[
 \operatorname{depth}_R N\le \operatorname{pd}_R M .
\]
This implies \(\dim_R N\le\operatorname{pd}_R M\) whenever
\(N\) is Cohen--Macaulay. The latter dimension bound also
holds for arbitrary \(N\) over a regular local ring.
\end{corollary}
\begin{proof}
Let \(F_\bullet\) be a finite free resolution of \(M\) of
length \(\operatorname{pd}_R M\), obtained by taking finite free
surjections through the final projective syzygy, which is finite
free over the local ring. All homology groups of
\(F_\bullet\otimes_R N\) have finite length by the localization
argument of \ref{res:I01}. Its zeroth homology \(M\otimes_R N\)
is nonzero by reducing modulo the maximal ideal and using
Nakayama. Now \ref{res:I14} gives the depth inequality.
When \(N\) is Cohen--Macaulay its depth equals its dimension.
When \(R\) is regular, \ref{res:I13} and the
Auslander--Buchsbaum formula give
\[
 \dim N\le\dim R-\dim M
          \le\dim R-\operatorname{depth}_R M
          =\operatorname{pd}_R M.
\]
The depth-dimension bound and the Auslander--Buchsbaum formula
are \cite[DP-02 and DP-13]{core}.
\end{proof}

\begin{example}[Tangency and excess Tor]\label{res:I16}
In \(k[[x,y]]\), the two quotients by \(y\) and by \(y-x^b\),
for \(b\ge1\), have intersection multiplicity \(b\). The
two axes in the plane have multiplicity one. In
\(k[[x,y,z]]\), the modules \(R/(x,y)\) and \(R/(y,z)\)
meet only at the closed point and have intersection
multiplicity zero even though their tensor product has length
one.
\end{example}
\begin{proof}
For the tangency, resolve \(R/(y)\) by the two free terms for
multiplication by \(y\). On \(R/(y-x^b)\cong k[[x]]\) this
becomes multiplication by \(x^b\), which is injective with
cokernel of length \(b\). For the two axes the same resolution
with the second quotient \(R/(x)\) leaves multiplication by
the variable \(y\) in \(k[[y]]\), with cokernel \(k\).
For the last pair resolve \(R/(x,y)\) by its Koszul
complex. On \(R/(y,z)\cong k[[x]]\), the first Koszul
factor on \(x\) has homology \(k\) only in degree zero,
and the second factor on \(y=0\) has differential zero.
The tensor complex therefore has \(k\) in degrees zero
and one and zero higher homology. This verifies the
cancellation directly and agrees with the dimension sum
\(1+1<3\).
\end{proof}

\begin{proposition}[The supported Euler pairing]\label{res:I17}
Let \(F_\bullet,G_\bullet\) be bounded complexes of finite
free modules over a Noetherian local ring. Suppose the unions
of the supports of their homology meet at most at the closed
point. The alternating length of the homology of
\(F_\bullet\otimes G_\bullet\) is finite and additive in
each variable on degreewise short exact sequences of such
bounded free complexes. A homology isomorphism between two
such bounded free complexes in one variable preserves the
pairing.
\end{proposition}
\begin{proof}
A bounded exact complex of finite free modules is contractible.
Starting from its lowest degree, its first surjective
differential splits because the target is free. The kernel is
then a direct summand of a free module and is projective.
Repeat the splitting at each succeeding degree. For each
nonmaximal prime, one factor of the localized tensor complex
is exact and hence contractible, so the total complex is
contractible there. Its homology is finite and supported at
most at the closed point, hence has finite length by
\ref{res:I01}'s length argument.

Tensoring a degreewise short exact sequence with a bounded free
complex produces a degreewise short exact sequence of totals.
The long exact homology sequence and length additivity imply
the assertion. The cone of a homology isomorphism between
bounded free complexes is a bounded exact free complex, so
it is contractible. Its tensor product is again contractible.
Because formation of the cone commutes with tensor product
up to the usual interchange sign, the induced tensor map
is a homology isomorphism. For two finite free resolutions
of modules the pairing recovers \ref{res:I01}. Thus the
same Euler expression makes sense on perfect complexes
with isolated intersection of supports.
\cite[Chapter 5, \S 5.3, and Chapter 11]{roberts}
\end{proof}

\section{General results in the literature}

The direct proofs above establish the dimension bound
in every characteristic, all signs in equal and unramified
mixed characteristic, and the complete-intersection case without
a characteristic restriction. The following statements describe
the broader literature and supply context for the chapter.

Roberts and Gillet--Soul\'e proved vanishing when
\(\dim M+\dim N<\dim R\) for general regular local rings.
The supported pairing of \ref{res:I17} is the algebraic
quantity to which the localized Chern character and
Adams operations are applied in those proofs. The local
Chern character associates a cycle operation to a
complex equipped with its support, rather than to its
ordinary rank. Tensor product then corresponds to composition
of the localized operations. The construction and its
Riemann--Roch formula are developed in
\cite[Chapter 11 and \S 13.1]{roberts}.
Gabber's nonnegativity theorem applies also in the ramified
mixed-characteristic case.
\href{https://doi.org/10.1090/S0273-0979-1985-15394-7}{Roberts's announcement}.
\href{https://numdam.org/item/SB_1995-1996__38__273_0.pdf}{Berthelot,
\S 6.1, Theorem 6.1.1}.

The September 2026
\href{https://github.com/openai/math/blob/main/preprints/Positivity-of-Serres-Intersection-Multiplicity-September-23-2026/paper.pdf}{preprint
\emph{Positivity of Serre's Intersection Multiplicity}}
reports strict positivity for complementary dimensions
also in the ramified mixed-characteristic case. The classical
direct argument in \ref{res:I12} uses the unramified
hypothesis precisely in its coefficient-DVR and power series
description. The dimension argument in \ref{res:I13}
uses a larger unramified ring and reads a dimension bound
from its Euler cancellation. That cancellation concerns
\(\chi_S\) over the larger ring and gives no formula for
\(\chi_R\) in the ramified ring.

The new intersection theorem concerns the length of a
nonexact bounded finite free complex with finite-length
homology over a general Noetherian local ring. It replaces
the depth bound of \ref{res:I14} by the Krull dimension
bound. This strengthening matters over non-Cohen--Macaulay
rings. The equicharacteristic setting is discussed in
\cite[\S 9.4]{bruns}. Roberts's mixed-characteristic proof
uses localized Chern characters. A proof using Adams
operations on \(K\)-theory with supports is described in
\href{https://arxiv.org/abs/0706.0201}{Piepmeyer and Walker,
\emph{A New Proof of the New Intersection Theorem}}.
